% ============================================================
\section{4.4 GMM, EM 알고리즘, 그리고 LDA}
\begin{frame}{닭이 먼저냐 달걀이 먼저냐: 잠재변수 문제}
  4.3절 k-means는 각 점을 \textbf{딱 하나의} 클러스터에 배정
  (\kb{하드 할당}). 실제 데이터는 경계가 \textbf{모호}한 경우가
  많다 --- 어떤 점은 ``A에 60\%, B에 40\%'' 속한다고 보는 게 더
  자연스러울 수 있다.
  \vskip0.8em
  \begin{itemize}
    \item \kb{가우시안 혼합모델(GMM)}: 각 클러스터를
      \textbf{정규분포}로 표현(Week 3 GDA와 같은 가정)하고,
      각 점이 어느 클러스터에서 왔는지를 \textbf{확률}로
      표현(소프트 할당)
    \item 문제는 각 점의 \textbf{진짜} 클러스터 $z$가 데이터에
      \textbf{관측되지 않음}(\kb{잠재변수}, latent variable)
  \end{itemize}
  \vskip0.6em
  \begin{block}{순환}
    \begin{itemize}\small
      \item $z$(진짜 클러스터)를 \textbf{안다면} $\to$ 각
        클러스터의 평균·분산 구하기 \textbf{쉬움}(GDA처럼)
      \item 평균·분산을 \textbf{안다면} $\to$ 각 점의
        \textbf{사후확률} $z$ 구하기 \textbf{쉬움}(베이즈 정리)
      \item \textbf{둘 다 모른다} = 닭 vs 달걀의 순환
    \end{itemize}
  \end{block}
  EM은 이 순환을 \textbf{반복}으로 풀어낸다: 하나를 (아무 값으로)
  고정 $\to$ 다른 하나를 구함 $\to$ 그걸로 다시 처음 것을 고침,
  \textbf{값이 더 이상 바뀌지 않을 때까지}.
\end{frame}
